% =========================
% appendix.tex
% =========================

\setcounter{equation}{0}
\renewcommand{\theequation}{\Alph{section}.\arabic{equation}}
\setcounter{lemma}{0}
\renewcommand{\thelemma}{\Alph{section}.\arabic{lemma}}
\setcounter{proposition}{0}
\renewcommand{\theproposition}{\Alph{section}.\arabic{proposition}}

\section{Gaussian tail identities}
\label{app:tails}
Let $Z\sim\mathcal N(0,1)$ and $\varphi(z)$, $\Phi(z)$ denote its pdf/cdf, where $\varphi(z)=\frac{1}{\sqrt{2\pi}}e^{-z^2/2}$ and $\Phi(z)=\int_{-\infty}^z \varphi(s)\,ds$. For any $a\in\mathbb R$, we derive the following truncated moment formulas by integration by parts:
\begin{align}
\int_a^{\infty} z\,\varphi(z)\,dz &= \varphi(a),\label{eq:tail1}\\
\int_a^{\infty} z^2\,\varphi(z)\,dz &= a\,\varphi(a) + 1-\Phi(a),\label{eq:tail2}\\
\int_a^{\infty} z^3\,\varphi(z)\,dz &= (a^2+2)\,\varphi(a),\label{eq:tail3}\\
\int_a^{\infty} z^4\,\varphi(z)\,dz &= (a^3+3a)\,\varphi(a) + 3\,(1-\Phi(a)).\label{eq:tail4}
\end{align}

\subsection*{A.1\quad Derivation sketch}
These identities follow from integration by parts using $\varphi'(z)=-z\varphi(z)$. For \eqref{eq:tail1}: $\int_a^\infty z\varphi(z)dz = -\varphi(z)\big|_a^\infty = \varphi(a)$. For \eqref{eq:tail2}: write $z^2\varphi(z)=-z\cdot\varphi'(z)$ and integrate by parts to get $-z\varphi(z)\big|_a^\infty + \int_a^\infty\varphi(z)dz = a\varphi(a) + [1-\Phi(a)]$. Higher moments follow recursively.

\subsection*{A.2\quad Application to our model}
Setting $a=-S$ and defining $X:=S+Z$, note that $\{X\ge 0\} = \{Z \ge -S\}$. Therefore:
\begin{align}
\Pr[Z\ge -S] &= \Phi(S),\\
\E[Z\,\mathbf 1_{\{Z\ge -S\}}] &= \varphi(S),\\
\E[Z^2\,\mathbf 1_{\{Z\ge -S\}}] &= \Phi(S) - S\,\varphi(S),\\
\E[X^2\,\mathbf 1_{\{X\ge 0\}}] &= (S^2+1)\,\Phi(S) + S\,\varphi(S),\\
\E[X^3\,\mathbf 1_{\{X\ge 0\}}] &= (S^3+3S)\,\Phi(S) + (S^2+2)\,\varphi(S),\\
\E[X^4\,\mathbf 1_{\{X\ge 0\}}] &= (S^4+6S^2+3)\,\Phi(S) + (S^3+5S)\,\varphi(S).
\end{align}

\begin{lemma}[Useful tail moments with threshold $-S$]
\label{lem:tails}
With $Z\sim\Normal(0,1)$, $X=S+Z$, and the indicator $\mathbf 1_{\{\cdot\}}$:
\begin{align*}
\Pr[Z\ge -S] &= \Phi(S), & \E[Z\,\mathbf 1_{\{Z\ge -S\}}] &= \varphi(S),\\
\E[Z^2\,\mathbf 1_{\{Z\ge -S\}}] &= \Phi(S) - S\,\varphi(S), & \E[X^2\,\mathbf 1_{\{X\ge 0\}}] &= (S^2+1)\,\Phi(S) + S\,\varphi(S).
\end{align*}
\end{lemma}

\section{AUM moments, variance, and derivatives}
\label{app:variance}
Let $\Kcal:=A_0+f_1$, $X=S+Z$, $I:=\mathbf 1_{\{X\ge 0\}}$. Recall that fund returns are $r=r_f+\sigma X$ and excess returns are $r-r_f=\sigma X$. The flow function is $F(r-r_f)=f_1(r-r_f)+f_2(r-r_f)^2 I$. Therefore, end-of-period AUM is
\begin{align}
A &= A_0(1+r_f) + \Kcal\,\sigma X + f_2\,\sigma^2 X^2 I.
\end{align}
The constant term $A_0(1+r_f)$ represents the riskless growth of initial AUM. The term $\Kcal\,\sigma X$ captures both (i) the risky return on initial AUM and (ii) the linear flow response to performance. The term $f_2\,\sigma^2 X^2 I$ represents the convex flow kicker that activates when excess returns are positive.

\subsection*{B.1\quad Mean of AUM}
Using $\E[X]=S$ and Appendix~\ref{app:tails} for $\E[X^2 I]$, we obtain
\begin{align}
\E[A] &= A_0(1+r_f) + \Kcal\,S\,\sigma + f_2\,\sigma^2\,C_1,\label{eq:EA-correct}\\
C_1 &:= (S^2+1)\,\Phi(S) + S\,\varphi(S).\label{eq:C1-app}
\end{align}

\subsection*{B.2\quad Variance of AUM}
Since the constant term $A_0(1+r_f)$ has zero variance, and writing $A - A_0(1+r_f) = \Kcal\,\sigma X + f_2\,\sigma^2 X^2 I$, we have
\[
\V[A] = \V[\Kcal\,\sigma X + f_2\,\sigma^2 X^2 I].
\]
Using the variance formula $\V[aU+bV]=a^2\V[U]+b^2\V[V]+2ab\,\Cov(U,V)$ with $U=X$ and $V=X^2 I$, and the fact that $\V[X]=1$ since $X=S+Z$ with $Z\sim\Normal(0,1)$, we obtain
\begin{align}
\V[A] &= \Kcal^2\,\sigma^2 + 2\,\Kcal f_2\,\sigma^3\,\Cov(X, X^2 I) + f_2^2\,\sigma^4\,\V(X^2 I).\label{eq:VarA-app}
\end{align}
Compute the cross and quadratic terms from Appendix~\ref{app:tails}:
\begin{align}
\E[X^2 I] &= C_1,\quad \E[X^3 I] = (S^3+3S)\Phi(S) + (S^2+2)\varphi(S),\label{eq:X3}\\
\E[X^4 I] &= (S^4+6S^2+3)\Phi(S) + (S^3+5S)\varphi(S).\label{eq:X4}
\end{align}
Hence, using the definition $\Cov(U,V)=\E[UV]-\E[U]\E[V]$ and noting that $\E[X]=S$:
\begin{align}
\Cov(X, X^2 I) &= \E[X^3 I] - S\,\E[X^2 I] \nonumber\\
&= (S^3+3S)\Phi(S) + (S^2+2)\varphi(S) - S\,C_1.\label{eq:Cov}
\end{align}
Similarly, for the variance of $X^2 I$:
\begin{align}
\V(X^2 I) &= \E[X^4 I] - C_1^2 \nonumber\\
&= (S^4+6S^2+3)\Phi(S) + (S^3+5S)\varphi(S) - C_1^2.\label{eq:VarX2I}
\end{align}

\subsection*{B.3\quad Derivatives in $\sigma$}
Differentiate with respect to $\sigma$:
\begin{align}
\partial_\sigma \E[A] &= \Kcal\,S + 2 f_2\,\sigma\,C_1,\label{eq:dEA-app}\\
\partial_\sigma \V[A] &= 2\Kcal^2\sigma + 6\Kcal f_2\,\sigma^2\,\Cov(X, X^2 I) + 4 f_2^2\,\sigma^3\,\V(X^2 I).\label{eq:dVarA-app}
\end{align}
Define
\begin{equation}
\Delta_{\text{quad}}(S,f_2;\sigma):=6\Kcal f_2\,\sigma^2\,\Cov(X, X^2 I) + 4 f_2^2\,\sigma^3\,\V(X^2 I),\label{eq:DeltaQuad}
\end{equation}
so that $\partial_\sigma \V[A]=2\Kcal^2\sigma+\Delta_{\text{quad}}(S,f_2;\sigma)$.

\section{Manager's problem, best response, and IC fee}
\label{app:manager-solution}

In this appendix we derive the manager's first-order condition and the incentive-compatible fee schedule used in Section~\ref{subsec:IC}. The appendix shows how the manager's optimal volatility choice depends on the marginal effect of volatility on expected AUM and AUM risk, and then derives the fee that implements a target volatility chosen by the investor.

\subsection*{C.1\quad Manager's First-Order Condition}
The manager's objective is to maximize mean-variance utility over fee revenue $\phi A$:
\[
\max_\sigma \; U_M(\sigma,\phi) = \phi\,\E[A] - \eta\phi^2\,\V[A],
\]
where $\eta>0$ is the manager's risk aversion coefficient. Taking the first-order condition with respect to $\sigma$ and using the derivatives from \eqref{eq:dEA-app} and \eqref{eq:dVarA-app}:
\begin{align}
\frac{\partial U_M}{\partial \sigma} &= \phi\,\partial_\sigma \E[A] - \eta\phi^2\,\partial_\sigma \V[A] \nonumber\\
&= \phi\Big(\Kcal S + 2 f_2\,\sigma\,C_1\Big) - \eta\phi^2\Big(2\Kcal^2\sigma + \Delta_{\text{quad}}\Big) \nonumber\\
&= 0.\label{eq:manager-FOC-app}
\end{align}
The left side represents the marginal net benefit of increasing volatility: the first term captures the increase in expected fee revenue (through both higher expected returns when $S>0$ and convex flow effects when $f_2>0$), while the second term captures the disutility associated with increased fee-revenue risk.

\subsection*{C.2\quad Existence, Uniqueness, and Second-Order Condition}
Under the condition $\eta\phi\cdot 2\Kcal^2 > 2f_2 C_1$, the marginal cost grows faster than the marginal benefit for large values of $\sigma$. As a result, the term $-\eta\phi^2\cdot 2\Kcal^2\sigma$ dominates for large $\sigma$, making \eqref{eq:manager-FOC-app} strictly decreasing. Since the FOC is positive at $\sigma=0$ (when $S>0$) and eventually negative for large $\sigma$, the intermediate value theorem guarantees a unique interior solution $\sigma^*(\phi)>0$.

\textit{Sufficiency condition.} To verify that the IC fee \eqref{eq:phi-IC-app} satisfies the above condition, substitute $\phi = \phi(\sigma_d)$ from \eqref{eq:phi-IC-app} into $\eta\phi\cdot 2\Kcal^2 > 2f_2 C_1$. In practice, this condition is verified numerically across the calibrated parameter ranges in Appendix~\ref{app:numerical}; it holds whenever $\Kcal\,S > 0$, which is ensured by $S>0$ and $A_0, f_1 > 0$.

\textit{Second-order condition.} The manager's objective $U_M(\sigma,\phi) = \phi\E[A] - \eta\phi^2\V[A]$ has second derivative
\[
\frac{\partial^2 U_M}{\partial\sigma^2} = \phi\,\partial_\sigma^2\E[A] - \eta\phi^2\,\partial_\sigma^2\V[A].
\]
Since $\partial_\sigma^2\E[A] = 2f_2 C_1 > 0$ (mean is convex in $\sigma$ due to flow convexity) and $\partial_\sigma^2\V[A] > 0$ (variance is convex in $\sigma$), the SOC $\partial^2 U_M/\partial\sigma^2 < 0$ requires the variance term to dominate: $\eta\phi\cdot 2f_2 C_1 < \eta\phi^2\,\partial_\sigma^2\V[A]$. This is equivalent to the IC condition above and is verified numerically over the calibrated parameter space. The second-order condition holds throughout the relevant range, confirming that the IC fee pins down a maximum rather than a minimum of the manager's objective.

\subsection*{C.3\quad Closed-Form Approximation}
To obtain additional analytical insight, it is useful to consider an approximation to the manager's best response. Under the linearized variance slope $\partial_\sigma \V[A]\approx 2\Kcal^2\sigma$, which captures the leading-order contribution of the linear term in AUM variance, the first-order condition simplifies. Solving for $\sigma$, the manager's best response admits the closed-form approximation
\begin{equation}
\label{eq:sigma-star}
\sigma^*(\phi) = \frac{\Kcal S}{2\,(\eta\phi\,\Kcal^2 - f_2 C_1)},\qquad \text{valid when } \eta\phi\,\Kcal^2> f_2 C_1.
\end{equation}

\subsection*{C.4\quad Incentive Compatibility}
Incentive compatibility requires that the manager's first-order condition holds at the target volatility $\sigma=\sigma_d$. Imposing this condition yields
\begin{equation}
\label{eq:phi-IC-app}
\phi(\sigma_d)
\;=\;
\frac{\partial_\sigma \E[A]\big|_{\sigma_d}}
{\eta\,\partial_\sigma \V[A]\big|_{\sigma_d}}
\;=\;
\frac{\Kcal\,S + 2 f_2 \,\sigma_d\,C_1}
{\eta\Big(2 \Kcal^2 \,\sigma_d + \Delta_{\text{quad}}(S,f_2;\sigma_d)\Big)}.
\end{equation}
This expression characterizes the fee schedule that implements the desired volatility target $\sigma_d$ by equating the manager's marginal benefit and marginal cost of risk-taking.

\section{Investor's FOC and equilibrium condition}
\label{app:investor}

\subsection*{D.1\quad Investor's First-Order Condition}
With net returns $r^{\text{net}}=(1-\phi)r-\phi$ and investor utility
\[
U_I(\sigma,\phi)=\E[r^{\text{net}}] - \frac{\gamma}{2}\V[r^{\text{net}}] = (1-\phi)(r_f+S\sigma)-\phi-\frac{\gamma}{2}(1-\phi)^2\sigma^2,
\]
the investor chooses the target volatility $\sigma_d$ to maximize $U_I(\sigma_d,\phi(\sigma_d))$, where the fee $\phi(\sigma_d)$ is pinned down by the incentive compatibility constraint \eqref{eq:phi-IC-app}.

Differentiating the investor's objective with respect to $\sigma_d$ using the chain rule:
\begin{align}
\frac{d U_I}{d \sigma_d} &= \frac{\partial U_I}{\partial \sigma}\Big|_{\sigma=\sigma_d} + \frac{\partial U_I}{\partial \phi}\Big|_{\phi=\phi(\sigma_d)} \cdot \frac{d\phi}{d\sigma_d} \nonumber\\
&= \Big[(1-\phi)S - \gamma(1-\phi)^2\sigma_d\Big] + \Big[-(r_f+S\sigma_d) -1 + \gamma(1-\phi)\sigma_d^2\Big]\frac{d\phi}{d\sigma_d}.\label{eq:investor-FOC-app}
\end{align}
Setting this equal to zero gives the equilibrium condition. The first bracketed term is the \textbf{direct effect}: increasing volatility raises expected net returns (by $S$) but also increases risk. The second term is the \textbf{contract adjustment}: changing the target requires adjusting the fee via the IC constraint, which affects investor utility through both the mean and the variance.

\subsection*{D.2\quad Equilibrium Characterization}
For notational convenience, define
\begin{align}
N(\sigma) &:= \Kcal\,S \;+\; 2 f_2 \sigma\,C_1, \\
D(\sigma) &:= \eta\Big(2 \Kcal^2 \sigma + \Delta_{\text{quad}}(S,f_2;\sigma)\Big), \\
\phi(\sigma) &:= \frac{N(\sigma)}{D(\sigma)},\qquad
\phi'(\sigma) \;=\; \frac{N'(\sigma)D(\sigma) - N(\sigma)D'(\sigma)}{D(\sigma)^2},
\end{align}
where $N(\sigma)$ is the numerator and $D(\sigma)$ is the denominator of the IC fee \eqref{eq:phi-IC-app}.

Plugging $\phi(\sigma)$ and $\phi'(\sigma)$ into \eqref{eq:investor-FOC-app} and defining
\begin{equation}
G(\sigma) := (1-\phi(\sigma))S - \gamma(1-\phi(\sigma))^2\sigma + \Big[-(r_f+S\sigma) -1 + \gamma(1-\phi(\sigma))\sigma^2\Big]\phi'(\sigma),
\end{equation}
the equilibrium target $\sigma_d^*$ satisfies $G(\sigma_d^*)=0$. This is a single nonlinear equation in one unknown that can be solved numerically via bisection or Newton's method.

\subsection*{D.3\quad Linear Approximation of the IC Fee}
For practical implementation, we linearize the IC fee schedule around a reference volatility $\sigma_0$ (typically chosen as the optimal target $\sigma_d^*$). Using a first-order Taylor approximation:
\begin{equation}
\phi(\sigma_d) \;\approx\; \phi(\sigma_0) + \phi'(\sigma_0)(\sigma_d - \sigma_0) = \underbrace{[\phi(\sigma_0) - \phi'(\sigma_0)\sigma_0]}_{\alpha} + \underbrace{\phi'(\sigma_0)}_{\beta}\,\sigma_d,
\end{equation}
where
\begin{equation}
\alpha := \phi(\sigma_0) - \phi'(\sigma_0)\,\sigma_0, \qquad \beta := \phi'(\sigma_0).
\end{equation}
The parameter $\alpha$ acts as a base fee, while $\beta$ captures the marginal fee adjustment with respect to the volatility target. This affine representation provides a transparent and practical implementation of the contract: once a volatility target is specified, the corresponding fee follows directly from a linear rule. As shown in Figure~\ref{fig:affine-approx}, this approximation is highly accurate in the neighborhood of the optimal target, with small approximation errors over empirically relevant ranges.

\section{Fee base equivalence}
\label{app:fee-base}
Different fee conventions can be reconciled as follows. Suppose fees are charged on beginning-of-period AUM at rate $\tilde\phi$, so the investor pays $\tilde\phi A_0$ and receives terminal wealth $A_0(1+r)$, giving a net return per dollar invested of
\[
r^{\text{net}} = \frac{A_0(1+r) - \tilde\phi A_0}{A_0} = r - \tilde\phi.
\]
In the model, fees are charged on end-of-period AUM at rate $\phi$, so the investor receives $(1-\phi)A$ from terminal AUM $A=A_0(1+r)$, giving a net return of
\[
r^{\text{net}} = \frac{(1-\phi)A_0(1+r)}{A_0} = (1-\phi)(1+r) - 1 = (1-\phi)r - \phi + O(\phi^2).
\]
Matching expected net returns to first order in fees:
\[
\E[r] - \tilde\phi \approx \E[r] - \phi(1+\E[r]) \quad\Rightarrow\quad \tilde\phi \approx \phi(1+r_f + S\sigma).
\]
For typical parameter values (e.g., $r_f=2\%$, $S\sigma \approx 5\%-10\%$, $\phi \approx 1\%-2\%$), the difference between conventions is second-order and negligible for practical purposes. Throughout the paper, we use the end-of-period convention as it aligns naturally with the flow-based AUM dynamics.

\section{Performance fees with high--water marks (sketch)}
\label{app:perf-fee}
Many hedge funds charge performance fees with high-water marks (HWM). A stylized HWM structure charges a performance fee $\psi(r-r_f)^+$ on returns exceeding the risk-free rate (or previous peak). With fund return $r=r_f+\sigma X$ where $X=S+Z$, this becomes
\[
\text{Performance fee} = \psi\cdot\max(r-r_f, 0) = \psi\sigma\cdot\max(X, 0) = \psi\sigma X^+.
\]
The manager's expected fee revenue now includes this performance component. Using $X^+ = X\cdot\mathbf{1}_{\{X\ge 0\}}$ and results from Appendix~\ref{app:tails}:
\[
\E[X^+] = \E[X\cdot\mathbf{1}_{\{X\ge 0\}}] = S\Phi(S) + \varphi(S),
\]
and higher moments follow similarly. Performance fees create additional convexity in the manager's payoff, similar to the flow convexity $f_2$ but with different functional form. The key difference is that HWM fees are capped (cannot exceed the actual gain), whereas flow convexity can amplify AUM growth unboundedly in extreme positive scenarios.

Incorporating both management and performance fees into the framework requires computing $\E[A]$ and $\V[A]$ including the performance fee impact on net returns and flows. The IC condition becomes more complex but follows the same economic logic: the fee structure (now $\phi$ and $\psi$) must be chosen to implement the desired volatility target. In practice, the equilibrium is computed numerically using the same bisection approach as in Appendix~\ref{app:investor}, with the added complexity that the manager's FOC now involves both fee components.

\section{Numerical implementation notes}
\label{app:numerical}

\subsection*{G.1\quad Solving for equilibrium}
To find the optimal target $\sigma_d^*$, solve $G(\sigma)=0$ where $G$ is defined in Appendix~\ref{app:investor}. We recommend bisection on an interval $[\underline\sigma,\overline\sigma]$ where:
\begin{itemize}
\item $\underline\sigma > 0$ is chosen small (e.g., $0.01$ or 1\% annualized) to avoid numerical issues near zero.
\item $\overline\sigma$ is chosen as an upper bound where the investor's direct utility gain from volatility becomes negative. A conservative choice is to set $\overline\sigma$ such that $(1-\phi)S - \gamma(1-\phi)^2\overline\sigma < 0$ when ignoring the contract adjustment term. For instance, with $\phi\approx 0.014$, $S=0.35$, $\gamma=5$, this gives $\overline\sigma \approx 0.35/(5\cdot 0.986) \approx 0.071$. In practice, we use $\overline\sigma = 0.40$ to be conservative.
\end{itemize}
Bisection is robust and converges rapidly since $G$ is continuous and typically has a unique zero in the relevant range.

\subsection*{G.2\quad Precomputing tail blocks}
For computational efficiency, precompute the Gaussian tail moments $\Phi(S)$, $\varphi(S)$, $C_1$, $\E[X^3 I]$, and $\E[X^4 I]$ once for each Sharpe ratio $S$ using the formulas in Appendix~\ref{app:tails}. These quantities depend only on $S$ (not on $\sigma$, $f_2$, etc.) and can be reused across all evaluations of $\phi(\sigma)$ and $G(\sigma)$ for that value of $S$.

\subsection*{G.3\quad Parameter ranges}
The baseline parameter values for the numerical exercise are $S=0.35$, $r_f=0.037$, $\gamma=5$, $\eta=3$, $A_0=1$, $f_1=1.5$, $f_2=25$. Comparative statics sweep the following ranges:
\begin{itemize}
\item Convexity $f_2 \in [5, 60]$
\item Initial AUM $A_0 \in [0.5, 3.0]$
\item Linear flow $f_1 \in [0.5, 3.0]$
\item Sharpe ratio $S \in [0.1, 1.0]$
\item Investor risk aversion $\gamma \in [2, 10]$
\item Manager risk aversion $\eta \in [1, 8]$
\end{itemize}
These ranges cover empirically plausible scenarios while ensuring the FOCs remain well-behaved. The comparative statics figures in Section~\ref{sec:numerical} are computed at the Section~\ref{sec:calibration} baseline reported above, sweeping each parameter over the ranges listed here while holding all others fixed. The key qualitative results --- in particular, the signs of $d\sigma_d^*/df_2$, $d\sigma_d^*/d\gamma$, $d\sigma_d^*/dS$, and the slope of the IC fee schedule --- hold robustly across all empirically plausible parameter combinations.

\subsection*{G.4\quad Affine approximation}
To construct the linear approximation $\phi(\sigma_d) \approx \alpha + \beta\sigma_d$, first solve for the optimal target $\sigma_d^*$, then compute $\phi(\sigma_d^*)$ and $\phi'(\sigma_d^*)$ using the formulas in Appendix~\ref{app:manager-solution} and Appendix~\ref{app:investor}. The approximation is then $\alpha = \phi(\sigma_d^*) - \beta\sigma_d^*$ and $\beta = \phi'(\sigma_d^*)$. As Figure~\ref{fig:affine-approx} demonstrates, this approximation remains accurate within $\pm 5\%$ of $\sigma_d^*$ with average errors around 0.1\%.
